\section{Introduction}
\label{sec:intro}


Self-improving agents aim to turn local improvements in coding and reasoning into a recursive loop, where each stronger agent can produce better variants of itself. Search methods such as the Darwin and Huxley G\"odel Machines~\citep{zhang2025darwin,wang2025huxley} recently set the open-source state-of-the-art on agentic coding tasks by editing their own code and retaining variants that improve an external utility signal. \hyperagents~\citep{zhang2026hyperagents} then extended this self-improvement beyond coding to task agents in arbitrary domains. Despite this progress, current search methods still depend on stationary evaluation fixed outside the improvement loop, unlike biological evolution, where each species adapts to competitors that evolve in turn. This dependence also constrains self-improvement in three settings: when a target task has no direct benchmark, for example paper writing and proof writing lack one while paper reviewing and proof grading do not~\citep{DBLP:journals/corr/abs-2408-06292,yamada2025ai,luong2025towards}; when evaluation is slow or weakly informative~\citep{levine2018learning,burger2020mobile,NIPS2017_453fadbd}; and when static benchmarks saturate or become vulnerable to reward hacking as agents improve~\citep{DBLP:journals/corr/abs-2603-25681,DBLP:journals/corr/abs-2203-04592}.

Addressing these concerns is a prerequisite for extending self-improvement to open-ended settings like scientific research and writing. We therefore introduce the \textbf{Red Queen G\"odel Machine (\ourmethod)}, an evolutionary framework that treats evaluation as part of the search process. We study this through \emph{co-evolved} learned evaluators, which improve alongside the task agents they guide, defining the search utility. Such an evaluator can supply the search signal when no benchmark exists, add criteria a benchmark cannot capture, like code maintainability, or provide a cheaper evaluation proxy.

The core mechanism enabling this co-evolution is \emph{controlled utility evolution}, which divides search into \emph{evolutionary epochs}. Within an epoch, one evaluator is frozen and grades every task agent, supplying a stationary utility signal; the utility can change only at epoch boundaries. In parallel, the search co-evolves challenger evaluators in a codebase shared with their task agents and scores them against a held-out ground-truth dataset.  At an epoch boundary, a challenger replaces the frozen evaluator only if it statistically outperforms the incumbent on this ground truth; \theourmethod then applies \emph{selective erasure}, discarding only utility records that depended on the replaced evaluator. Because each epoch is therefore a \emph{fixed-criterion search problem}, prior self-improvement guarantees apply directly~(\cref{sec:method,app:proofs}), while the search objective can evolve across epochs, for instance by promoting a stronger evaluator or adding an adversarial-sample regularizer.


\begin{figure}[t]\centering\includegraphics[width=\linewidth]{figures/fig1_polyglot_arrows.pdf}

  \caption{\textbf{\ourmethod\ exceeds the prior SOTA \hgmh\ on \polyglot with $1.35\times$--$1.72\times$ fewer search tokens by adding a cheap evolved code reviewer.} \emph{Left:} held-out pass rate vs.\ search cost; arrows give the tokens \ourmethod\ saves to exceed the baseline's rate, and the heart marks the run's best agent. \emph{Right:} best-belief utility during search. At each evaluator replacement (crowned dashed rule), the utility drops as selective erasure discards records scored by the displaced reviewer, then re-climbs under its replacement; shaded bands mark epochs. Full details in \cref{sec:experimental-design,exp_design:reading_results_figures}.}
    \label{fig:intro-token-efficiency}
\end{figure}

We provide a preliminary empirical investigation of \theourmethod across three domains following \citet{zhang2026hyperagents}: coding~\citep{gauthier2024polyglot}; paper writing and reviewing~\citep{lu2024ai,zhao2026apres}; and proof writing and grading~\citep{zhang2026hyperagents,luong2025towards}. On verifiable coding tasks, \theourmethod reaches a held-out pass rate of $\textbf{71.7}\boldsymbol{\%}$ against the prior SOTA's $69.9\boldsymbol{\%}$~(\cref{fig:intro-token-efficiency}) by adding an agent-as-a-judge code reviewer. Its quality signal complements test execution, improving the pass rate, while requiring fewer tokens compared to multi-turn coding agent execution. This illustrates a broader advantage of co-evolution: shared expansions can improve multiple utilities at once~(see \cref{fig:analysis-transfer}). We then turn to scientific paper writing and reviewing, and Olympiad-level proof writing and grading, where \theourmethod outperforms prior self-improving agents: co-evolved writers raise the acceptance rate of their papers from the prior SOTA's $21.8\boldsymbol{\%}$ to $\textbf{40.5}\boldsymbol{\%}$~(\cref{tab:paper-writer-cross-reviewer}), while a co-evolved grader exceeds the static baselines with a $\textbf{3}\boldsymbol{\times}$ lower search cost than the prior SOTA~(\cref{fig:prover-grader-results}). We further observe that as evaluators strengthen across epochs, they may impose a curriculum-like effect on the task agents, with each transition re-ranking the archive~(\cref{fig:transition-rerank,fig:tree-paper}). Co-evolution can also regularize evaluators: the strongest baseline reviewer over-accepts \AI-generated papers relative to human ones~\citep{DBLP:conf/nips/PanicksseryBF24}, but \theourmethod's evolutionary epochs let us first gather \AI-generated samples that a frozen evaluator accepts, then replay them as adversarial samples in later epochs, producing a reviewer equally stringent on machine and human work~(\cref{fig:paper-writer-reviewer-results}).


\emph{We believe \theourmethod points toward a class of self-improving co-evolving systems in which agents and evaluators recursively bootstrap one another without further human intervention to reach capabilities beyond those of static evaluation . If successful, this could be a meaningful step toward more capable \AI{} systems, at the cost of loosening convergence guarantees compared to static evaluation criteria.}
\introtakeawaybox{
 \begin{enumerate}[leftmargin=*]
 \item \textbf{Recursive Self-Improvement under Non-Stationary Utilities.}
 \Theourmethod extends self-improving agents beyond fixed external evaluation by treating the utility signal as part of the search process. This allows search objectives to evolve, opening self-improvement to learned non-stationary utilities~(\cref{sec:method}).
 \item \textbf{Controlled Utility Evolution.}
 We make non-stationary utilities compatible with prior self-improvement guarantees via \emph{controlled utility evolution}: search runs in fixed-evaluator epochs, and the objective changes only at epoch boundaries. A challenger evaluator is promoted only when it raises an $\epsilon$-best-belief score on ground truth, a lower bound its utility exceeds with probability $1-\epsilon$~(\cref{sec:method,app:proofs}).
 \item \textbf{Practical Benefits of Co-Evolved Learned Evaluators and Search Interventions.}
 Empirically, co-evolution adds a cheap learned code-review signal that improves search efficiency even where a verifier exists~(\cref{sec:exp-polyglot}), enables generator search where direct benchmarks are unavailable~(\cref{sec:exp-writer}), and enables debiasing over-lenient evaluators through adversarial objectives~(\cref{sec:exp-adversarial}).
 \end{enumerate}
 }